\section{Results}
\label{sec:results}

\paragraph{The object.} One Qwen3-4B~\citep{qwen3_2025} file throughout:
3{,}633{,}315{,}840 weights in 252 projections, sealed as
\texttt{qwen3-4b-dclm-ft.bin}, 1{,}794{,}564{,}765~B. It holds 216 \Tetra{}
records, 36 int4~g128 records for \texttt{v\_proj}, the exactly-kept tail and
the row scales; the embedding is quantized to int8 at load. Calibration is
131{,}072 tokens of DCLM-edu~\citep{dclmedu}, followed by the row-scale
training of \S\ref{sec:rowscales}.

\paragraph{Two references, kept apart.} The \emph{f16 checkpoint} is the
unquantized model, 8.04~GB. The \emph{dense reconstruction} is this file
decoded back to f16 tensors and run through an ordinary forward pass. Every
perplexity and MMLU number below, ours and the baselines', is a dense
reconstruction; the throughput numbers are the served kernel.
\S\ref{sec:served} gives the numerical link between the two.

\paragraph{Protocol.} Kernel timings are one process on one L40S, all arms
interleaved in the same order, seven rounds with the first two discarded,
ratios formed round by round. They cover the 252 projection matrix--vector
products of one token, which is less than a token. Perplexity is
wikitext-2~\citep{wikitext2016} at context 4096 over 12 windows in f16, with a
token fingerprint that must match across arms. MMLU~\citep{mmlu2021} is 5-shot
micro-averaged over the full 14{,}042-question census. The 2{,}280-question
sample used in our earlier paper differs from the census by 0.18~pp on the f16
checkpoint and by 1.15~pp on a 2-bit object, so cells from the two are marked
and never subtracted.

\subsection{Bytes and time}

\begin{table}[t]
\centering\footnotesize\setlength{\tabcolsep}{4pt}
\begin{tabular}{lrrrrr}
\toprule
 & FP16 & AWQ w4g128 & \Planes{} & \Tetra{} & control\\
\midrule
matrices timed & 252 & 252 & 252 & 216 & 252\\
median ms & 10.973 & 3.261 & 4.997 & 3.424 & 2.289\\
GB read per pass & 7.27 & 1.90 & 2.18 & 0.95 & 0.07\\
\bw{} in VRAM & 16.000 & 4.179 & 4.804 & \textbf{2.148} & 0.159\\
GB/s & 662 & 583 & 437 & 278 & 32\\
\bottomrule
\end{tabular}
\caption{Kernel benchmark, tile 64, one process, one L40S. The control is
\texttt{nullk}, the same launch geometry reading no weight byte. The \Tetra{}
arm times the 216 lattice matrices of the object; the 36 int4 matrices are read
by a different kernel and are not timed, so its millisecond and GB cells cover
less work than the other arms'. Bits per weight is the row that compares. The
ten-arm table is Appendix~\ref{app:bench}. Data: echelle-formats.csv.}
\label{tab:bench}
\end{table}

Per weight, \Tetra{} reads 2.148 bits where the 4-bit baseline reads 4.179 and
\Planes{} reads 4.804: less than half the traffic of AWQ, and less than half
that of the layout it replaces.

The millisecond column does not rank the arms. The \Tetra{} arm covers 216 of
the 252 matrices while every other arm covers all 252, so its 3.424~ms and
AWQ's 3.261~ms are different amounts of work, and subtracting the control from
both does not fix that. The measurement that would rank them is in
\S\ref{sec:limitations}. What the timings do show is the size of the fixed cost
around the decode: the control, which reads no weights at all, accounts for
2.289~ms of arms that run between 3.4 and 5.0~ms.

\subsection{In the model}
\label{sec:served}

Roughly 48\,\% of a token is spent outside the projections, so the benchmark's
milliseconds do not convert into tokens per second. Measured in the model, with
the served flags, batch 1, greedy, five rounds, the object runs at
\textbf{98.3~tokens/s} [97.5--98.6] in \textbf{1.39~GB}, against 43.4
[43.4--43.4] in 8.04~GB for the f16 checkpoint in the same process:
$\times2.27$ in speed, $\div5.81$ in memory. The 1.39~GB is a host count of
weights, 0.97~GB at 2.140~\bw{} with the tail in binary16, plus 0.41~GB of int8
embedding; it excludes the KV cache and the activations. The run was taken at
$T=128$, before the tile of \S\ref{sec:tile} moved; throughput at $T=64$ has
not been measured.

\paragraph{What validates the kernel.} The per-row check does, not the tokens.
Before any arm is timed, each of the 1{,}105{,}920 output rows of the model is
recomputed in f64 from the same dense reconstruction and compared against the
kernel's row, at a threshold of $10^{-5}\cdot\sum|w\cdot x|$; the \Tetra{}
arm's worst row is $2.1\times10^{-8}\cdot\sum|w\cdot x|$. One level down, the
device decoder's 24 outputs per block are compared bit for bit against an
independent host implementation of the same word, on the development machine
and again on the card. The token check is weaker and end to end: the served
object's first 32 greedy tokens agree with the f16 checkpoint's on one prompt,
and an earlier \Tetra{} object agreed over 256. The \texttt{oracle} binary that
the journals record beside it compares our forward implementation with
\texttt{candle-transformers} on an unquantized checkpoint, so it says nothing
about the quantized kernel.

\begin{table}[t]
\centering\footnotesize\setlength{\tabcolsep}{4pt}
\begin{tabular}{llp{6.4cm}}
\toprule
count & value & what it includes\\
\midrule
on disk & 1.79~GB & all parameters, embedding included: 2.7475~b/param\\
read per pass & 0.95~GB & the 216 \Tetra{} matrices, 2.148~\bw{} (bench accounting, f32 tail)\\
resident, serving & 1.39~GB & 0.97 projections (2.140~\bw{}, f16 tail) $+$ 0.41 int8 embedding; no KV cache, no activations\\
DRAM traffic & --- & not measured; we report the bytes the kernel must read, not counters\\
\bottomrule
\end{tabular}
\caption{Four byte counts of one object. They answer different questions and
are never divided against each other. For the pure-\Tetra{} reference file the
kernel rate decomposes exactly: 1.9907 (code) $+$ 0.1494 (f32 tail over
16{,}957{,}440 kept weights) $+$ 0.0097 (1{,}105{,}920 row scales) $=$
2.1498.}
\label{tab:accounting}
\end{table}

\subsection{Quality}

\begin{table}[t]
\centering\footnotesize\setlength{\tabcolsep}{4pt}
\begin{tabular}{lrrrrr}
\toprule
 & disk (GB) & b/param & VRAM \bw{} & wiki ppl & MMLU\\
\midrule
FP16 checkpoint & 8.04 & 16.000 & 16.000 & \textbf{12.2361} & \textbf{70.14}\\
AWQ w4g128 & 2.67 & 5.302 & 4.179 & 13.5207 & 70.04$^\dagger$\\
QTIP 2-bit & --- & --- & 2.000 & 17.04$^\ddagger$ & 57.4$^\ddagger$\\
IQ2\_XXS (llama.cpp) & \textbf{1.26} & \textbf{2.4967} & --- & --- & 38.87$^\dagger$\\
\Tetra{}, trained scales & 1.79 & 2.7475 & \textbf{2.148} & 12.3268 & 61.11\\
\bottomrule
\end{tabular}
\caption{Qwen3-4B, five formats. Bits per parameter cover the whole model,
embedding included. $^\dagger$2{,}280-question sample. $^\ddagger$Measured in
another paper's harness, on another corpus and protocol; those two cells are
context, not a comparison. Data: tetra-formats.csv.}
\label{tab:formats}
\end{table}

On the census, the object scores 9.03~pp of MMLU below the f16 checkpoint at
1.007 times its perplexity. The 4-bit baseline is within 0.3~pp of the
checkpoint on its own sample. Against the 4B rows of \citet{llvq2026}, 61.11
sits above their 60.7 without fine-tuning and 1.69~pp below their 62.8 with
it; the second is the comparable row, since this object is fine-tuned too.
Their numbers come from their harness and their calibration, which is why we
do not turn the gap into a ranking.

\subsection{Training the row scales}
\label{sec:rowscales}

Equation~\eqref{eq:recon} ends with $\sigma_{\mathrm{row}}$, and the file
already stores 1{,}069{,}056 of them for the 216 lattice matrices, fitted by
least squares at encoding time. We retrain those and nothing else. The lattice
directions stay frozen: the decoder is a table lookup, and the trainer refuses
a gradient into the decoded directions. No parameter is added, no field widens,
the decoder is byte-identical afterwards and the bit rate is unchanged. The
run costs 1.71~h on one L40S, about \$4: KL against the f16 checkpoint as
teacher on DCLM-edu, sequences of 1024, batch 2, learning rate
$3\times10^{-4}$, 100 warmup steps, bf16, 9{,}507 steps, 19{,}470{,}336 tokens.
Both evaluation sets are disjoint from that corpus. The trained multipliers are
folded back into the row-scale field: 216 matrices scaled, 36 int4 records
untouched.

\begin{table}[t]
\centering\footnotesize\setlength{\tabcolsep}{5pt}
\begin{tabular}{lrrrl}
\toprule
base & MMLU before & after & $\Delta$ & McNemar\\
\midrule
bare \Tetra{} (2.1498~\bw{}) & 54.64 & 58.94 & $+4.30$ & $p=4.4\cdot10^{-40}$\\
served object (2.148~\bw{}) & 57.95 & 61.11 & $+3.15$ & $p=1.8\cdot10^{-25}$\\
\bottomrule
\end{tabular}
\caption{The same training applied to two starting points, scored question by
question on a fixed file: same codes, same codebook, same byte count.
Perplexity of the served object falls from 16.2509 to 12.3268. Data:
tetra-rowscales.csv.}
\label{tab:rowscales}
\end{table}

The weaker starting point gains more, the direction \citet{llvq2026} report
over five pairs in their Table~6; two runs agree with it and do not establish
it. On the served object the training moves perplexity by $-24.1\,\%$, to
within 0.74\,\% of the f16 checkpoint, and MMLU by $+3.15$~pp, to 9.03~pp below
it (Fig.~\ref{fig:dissoc}). That the two metrics move so differently is
expected: the objective was KL against the checkpoint on text, which is close
to what perplexity measures, while MMLU was not optimised. The same split
appears, larger, in the fine-tuned 2-bit rows of \citet{llvq2026}, which read
below their own f16 baseline in perplexity while staying 7.4 points under it in
MMLU. On an object trained this way, perplexity alone is not enough to judge
quality.

\begin{figure}[t]
\centering
\includegraphics[width=0.60\linewidth]{fig_dissoc.pdf}
\Description{Scatter of MMLU against perplexity ratio. The FP16 checkpoint sits
at ratio 1.0 and 70.14 MMLU; the object before training at 1.33 and 57.95;
after training at 1.007 and 61.11, joined by an arrow.}
\caption{The served object before and after the row-scale training, against the
f16 checkpoint, on both metrics at once. Data: tetra-rowscales.csv.}
\label{fig:dissoc}
\end{figure}
